


Let  $F$ be a \pql map as in the previous subsection. The results in this section are purely combinatorial, so for simplicity may imagine that $F$ is a quadratic polynomial, so in particular $\iota$ is an inclusion (see \cite[Section 4.1]{K06} for details of transferring combinatorial data from \pql maps to quadratic polynomials).

We recall the definition of \emph{domination} from Appendix \ref{sec:domination}. Let us say that a WAD diagram on $\UU\sm \KK$ is \emph{horizontal} if it is supported on horizontal arcs. 
Let $\val$ be a function that assigns a non-negative real number to every WAD on 
$\UU^n\sm \KK^n$ with $n\geq 0$. If 
\begin{enumerate}
    \item $\val(X)>0$ for any horizontal WAD $X$;
    \item for any WADs $X$ on $\UU^n\sm \KK^n$ and $Y$ on $\UU\sm \KK$, if $X\dominates Y$ then $\val(X)\geq \val(Y)$; 
    \item there is some $\lambda>1$ so that $\val((f^*)^nX)\leq \lambda^{-1}\val(X)$ for every WAD $X$ on $\UU\sm \KK$ and every $n\geq 0$;
\end{enumerate}
then we will call $\val$ an \emph{$\lambda$-expanding valuation} \index{weighted arc diagram (WAD)!expanding valuation} for $F$. 
Recalling the core eigenvalue $\lambda_F$ from Section \ref{sec:hubbard trees}, we observe that a $\lambda_F$-expanding valuation always exists:



\begin{theorem}\label{thm:value}
    If $F$ is a primitively renormalizable and periodically repelling \pql-map with core eigenvalue $\lambda_F$, then there is a $\lambda_F$-expanding valuation for $F$.
\end{theorem}

Before proving Theorem \ref{thm:value}, we need the following two arithmetic facts:
\begin{lemma}
	Suppose that $a_i, b_i, c_i \ge 0$ for $i = 1 \ldots n$,
	and,
	for each $i$,
	\begin{equation} \label{eq:upper}
		a_i^2 + b_i^2 \le c_i^2.
	\end{equation}
	Then 
	\begin{equation} \label{eq:combined}
		\left(\sum_{i = 1}^n a_i\right)^2 + \left(\sum_{i = 1}^n b_i\right)^2 \le \left(\sum_{i = 1}^n c_i\right)^2.
	\end{equation}
\end{lemma}
\begin{proof}
	By Cauchy-Schwarz and \eqref{eq:upper},
	for all $i$ and $j$,
	\begin{equation} \label{eq:cs}
		a_i a_j + b_i b_j \le \sqrt{(a_i^2 + b_i^2)(a_j^2 + b_j^2)} \le c_i c_j.
	\end{equation}
	Adding this up over all $i$ and $j$, we obtain \eqref{eq:combined}.
\end{proof}

\begin{lemma}\label{lem:harmonic and geometric sum}
    For any sequences $(a_j)$ and $(w_j)$ of nonnegative real numbers, 
    \begin{equation}\label{eq:harmonic sum bound}
        \left(\bigoplus w_j\right)\left(\sum a_j\right)^2\leq \sum w_ja_j^2. 
    \end{equation}
\end{lemma}
\begin{proof}
    See \cite[Lemma 10.1]{K06}.
\end{proof}


\begin{proof}[Proof of Theorem \ref{thm:value}]
    We recall the linear  transformations $T_F: \R^{E(\TT)}\to \R^{E(\TT)}$ with eigenvalue $\lambda_F$ from Section \ref{sec:hubbard trees}. We have a natural bijection between $E(\TT)$ and $\HH^\biv$, so we can view any non-negative element of $\R^{E(\TT)}$ as a WAD with support in $\HH^\biv$. From this viewpoint, $T_F$ and its transpose $T_F^\top$ are exactly the linear transformation generated by 
    $$\beta\mapsto \sum_{\alpha\in \iota^*\beta} f_*\alpha \quad \text{ and } \quad \beta\mapsto \sum_{\alpha \in f^*\beta}\iota_*\alpha$$
    respectively.  {(Informally speaking, the former summation is taken over all $f$-push-forwards of $\beta$ (rel $\TT$),
    while the latter is taken over all $f$-pullbacks (rel $\TT$).}
    %using the convention that $f^*\beta = \setofst{\alpha}{f_*\alpha = \beta}$ and $\iota_*\alpha$
    %is the unique bivertical arc that pulls back to $\alpha$. 
    Additionally, 
    as $f^n\circ \iota_n^{-1}$ and $\iota_n\circ f^{-n}$ induce the linear transformations $T_F^n$ and $(T_F^\top)^n$ respectively, $f^n_*\iota_n^*\beta = (f_*\iota^*)^n\beta$ for all $\beta\in \HH^\biv$ and $n\geq 0$.


    Let $v$ and $w$ be WADs corresponding to eigenvectors with eigenvalue $\lambda_F$ of $T_F$ and $T_F^\top$ respectively and fix some $n\geq 0$. Thus 
    \begin{align*}
        \lambda_F^nv(\beta) =\sum_{\alpha\in \iota_n^*\beta}v(f^n_*\alpha) \text{ and }
        \lambda_F^nw(\beta) =\sum_{\alpha\in (f^n)^*\beta}w((\iota_n)_*\alpha)
    \end{align*}
    for all $\beta\in \HH^\biv$.
   {For any $\beta\in \HH^\biv$ and {$\alpha\in \iota_n^*(\HH^\biv)$} 
    %$\alpha\in \iota_n^*\beta$ 
    we set 
    \begin{equation*}
    \l(\beta)= v(\beta)\sqrt{w(\beta)},\quad 
    \l(\alpha) = \lambda_F^{-n} \, v(f_*^n\alpha)\sqrt{w((i_n)_*\alpha)}.
    \end{equation*}} 
    Thus 
\begin{align} \label{eq:basic-length}
	\sum_{\alpha \in \iota_n^*\beta} l(\alpha) = l(\beta)
\end{align}
and 
\begin{equation} \label{eq:basic-value}
	\sum_{{\alpha \in f_*^n \beta}} l(\alpha)^2 = \lambda_F^{-n} l(\beta)^2.
\end{equation}


For any $n\geq 0$ and arc $\gamma$ in $\UU^n\sm \KK^n$, we set $\l(\gamma) = \sum_{\alpha\in \iota_n^*\HH^\biv} \pair{\alpha}{\gamma}\l(\alpha)$.
For any arc $\beta$ in $\UU\sm \KK$, we observe that
\begin{equation} \label{eq:length}
	 \sum_{\alpha\in \iota_n^*\beta} l(\alpha) \ge l(\beta)
\end{equation}
and
\begin{equation} \label{eq:value}
	\sum_{\alpha\in (f^n)^*\beta} l(\alpha)^2 \le \lambda_F^{-n} l(\beta)^2. 
\end{equation}
Indeed, the inequality \eqref{eq:length} follows easily from \eqref{eq:basic-length},
and \eqref{eq:value} follows from \eqref{eq:basic-value} and the following:


Now, we define the {valuation} $\val$ by $$\val(X) = \sum_\alpha X(\alpha) l(\alpha)^2.$$
We then observe that,
because of \eqref{eq:length} and \eqref{eq:harmonic sum bound}, 
\begin{equation}
	X \dominates Y \implies \val(X) \ge \val(Y).
\end{equation}
Moreover, 
by \eqref{eq:value},
$\val((f^n)^*X) \le \lambda_F^{-n} \val(X)$.
\end{proof}

Moreover, our proof of Theorem \ref{thm:value} allows us to estimate the value of a WAD in terms of its size:
\begin{prop}\label{prop:value estimate}
    Let $v$ and $w$ be eigenweights of $T_F$ and $T_F^\top$ respectively.
    We can choose the valuation in Theorem \ref{thm:value} so that if $X$ is a WAD on $\UU\sm \KK$ and $\langle \alpha, \HH^\biv\rangle\leq C$ for all $\alpha\in \supp X,$
    then
    $$u_{\min}^2\|X\|\leq V(X)\leq C^2u_{\max}^2\|X\|,$$
    where $u_{\min}$ and $u_{\max}$ are the maximum and minimum values of $v(\alpha)\sqrt{w(\alpha)}$ over all arcs $\alpha\in \HH^\biv$ that intersect the support of $X$.
\end{prop}
    




%
%Let $F$ be
%critically periodic quadratic polynomial (or we can let it be any degree $ \ge 2$), 
%and we let $\DD_F$ be the closure of the immediate basin of the periodic cycle for $F$. 
%We let $H_F$ be the Hubbard tree for $F$.
%If $\alpha$ is an edge of $F\inv H_F$,
%we let $F(\alpha)$ be the edge of $H_F$ that it maps to,
%and $i(\alpha)$ be the edge of $H_F$ that it lies in.
%If $\alpha$ is an edge of $H_F$ or $F\inv H_F$,
%we let $\alpha^\perp$ be the ``bivertical dual'' of $\alpha$: the arc from $\infty$ to $\infty$ in $\C \sm \DD_F$ that crosses $\alpha$ exactly once, and crosses no other edge of $H_F$ (or $F\inv H_F$). 
%We let $H_F^\perp$ be the set of all $\alpha^\perp$ over all edges $\alpha$ in $H_F$; 
%we can then define $F, i\from F\inv H_F^\perp \to H_F^\perp$.  
%We can similarly apply $F$ and $i$ to the arcs in the ``bivertical dual'' to $H_F$:
%these are the arcs that connect $\infty$ to $\infty$, crossing a unique edge of $H_F$ exactly once. 
%
%By the Perron-Frobenius theorem,
%there exists a positive-real-valued function $v$ on $H_F$,
%such that 
%\begin{equation}
%\sum_{i(\alpha) = \beta} v(F(\alpha)) = \lambda v(\beta)
%\end{equation}
%for all edges $\beta$ of $H_F$. 
%Likewise there is $w$ on $H_F$,
%such that
%\begin{equation}
%\sum_{F(\alpha) = \beta} w(i(\alpha)) = \lambda w(\beta)
%\end{equation}
%for all $\beta$. 
%The value of $\lambda > 0$ is unique, and $w$ and $v$ are determined up to rescaling. 
%
%If we then let 
%\begin{equation}
%l(\alpha) = \lambda^{-1} v(F(\alpha)) \sqrt{w(i(\alpha))}
%\end{equation}
%for $\alpha \in F\inv H_F$
%and 
%\begin{equation}
%l(\beta) = v(\beta) \sqrt{w(\beta)}
%\end{equation}
%for $\beta \in H_F$,
%then 
%\begin{equation} \label{eq:basic-length}
%\sum_{i(\alpha) = \beta} l(\alpha) = l(\beta)
%\end{equation}
%and
%\begin{equation} \label{eq:basic-value}
%\sum_{F(\alpha) = \beta} l(\alpha)^2 = \lambda\inv l(\beta)^2.
%\end{equation}
%
%Now, 
%for any arc $\gamma$ on $\C \sm \DD_F$, 
%we let $l(\gamma) = \sum_{\alpha \in H_F} l(\alpha) \pair{\alpha^\perp}{\gamma}$.
%We likewise define $l(\gamma)$ for $\gamma \in \C \sm F\inv \DD_F$,
%using our definition of $l$ on $F\inv H_F$. 
%
%We then have that
%\begin{equation} \label{eq:length}
%(\alpha_i) \to \beta \implies \sum l(\alpha_i) \ge l(\beta)
%\end{equation}
%and, 
%for any arc $\beta$ in $\C \sm \DD_F$,
%\begin{equation} \label{eq:value}
%\sum_{F(\alpha) = \beta} l(\alpha)^2 \le \lambda\inv l(\beta)^2. 
%\end{equation}
%The inequality \eqref{eq:length} follows easily from \eqref{eq:basic-length},
%and \eqref{eq:value} follows from \eqref{eq:basic-value} and the following:
%\begin{lemma}
%Suppose that $a_i, b_i, c_i \ge 0$ for $i = 1 \ldots n$,
%and,
%for each $i$,
%\begin{equation} \label{eq:upper}
%a_i^2 + b_i^2 \le c_i^2.
%\end{equation}
%Then 
%\begin{equation} \label{eq:combined}
%\left(\sum_{i = 1}^n a_i\right)^2 + \left(\sum_{i = 1}^n b_i\right)^2 \le \left(\sum_{i = 1}^n c_i\right)^2.
%\end{equation}
%\end{lemma}
%\begin{proof}
%By Cauchy-Schwarz and \eqref{eq:upper},
%for all $i$ and $j$,
%\begin{equation} \label{eq:cs}
%a_i a_j + b_i b_j \le \sqrt{(a_i^2 + b_i^2)(a_j^2 + b_j^2)} \le c_i c_j.
%\end{equation}
%Adding this up over all $i$ and $j$, we obtain \eqref{eq:combined}.
%\end{proof}
%
%Now, for any horizontal weighted arc-diagram $Y$ on $\C \sm \DD$,
%we let $V(Y) = \sum_\alpha Y(\alpha) l(\alpha)^2$,
%and likewise for $X$ on $\C \sm F\inv \DD$. 
%We then observe that,
%because of \eqref{eq:length},
%\begin{equation}
%X \dominates Y \implies V(X) \ge V(Y).
%\end{equation}
%Moreover, 
%by \eqref{eq:value},
%$V(F^*X) \le \lambda\inv V(X)$.
%
%We can then conclude that there are no positive horizontal solutions to $F^* X \dominates X$. 
